\scanpage{385}\mspage{124}

e) J'ignore par contre si les flèches

\[
C_\infty(\underline t')\unsure{-|\longrightarrow}\deleted{\unsure{C}}C_\infty(t)
\]

sont des iso dans $\mathrm{Cat\,hot}$, bien qu'elles soient définies en termes d'hom. de systèmes inductifs de catégories, qui sont des iso. dans $\mathrm{Ind\,Cat\,hot}$. Il en est ainsi, même pour les flèches \deleted{\unsure{entre}}

\[
C_\infty(\underline t),C_\infty(\underline t')\rightrightarrows C_\infty(t\times t')
\]

utilisées pour comparer $C_\infty(t)$ et $C_\infty(t')$ canoniquement. \unsure{Ainsi|Aussi}, je ne vois pas comment définir des flèches en sens inverse dans $\mathrm{Cat}$, qui seraient candidates à être des inverses dans $\mathrm{Cat\,hot}$.

f) Dans le cas ``géométrique'' ($\unsure{\overline{\mathrm{di}}|\overline{\mathrm{Di}}}=\underline{\mathrm{Ty}}$, situation décrite en termes des $X,\overline\Sigma$, \deleted{\unsure{$a,b\in$}} et $a,b\in\mathrm{Ob}\,X$), d'où des $C_\infty(\underline t;a,b)$, $C_\infty(a,b)$, $H(a,b)$ etc, on n'a \uline{pas} \inserted{d'op.} de composition des chemins au niveau des \uline{catégories} $C_\infty(\underline t;a,b)$ et $C_\infty(t;b,c)$, mais seulement au niveau des objets \deleted{$C_\infty(a,b)$.} $H(a,b)$, $H(b,c)$ de $\mathrm{Ind\,Cat\,hot}$ ou de $\mathrm{Ind\,Hot}$\unsure{,|.} Suis-je seulement sûr de savoir définir \uline{dans Hot}

\[
\underline C_\infty(a,b)\times C_\infty(b,c)\longrightarrow C_\infty(a,c)\qquad ?
\]

Sûrement, il n'y a pas de \unsure{problèmes|problème} \deleted{pour} les images dans $\widehat{\mathrm{Hot}_0}$, puisque là ce sont les \unsure{$\simeq$|mêmes} que celles des $H(a,b)$, $H(b,c)$, $H(a,c)$ respectivement, et on a alors accouplements \unsure{parfaits}, associatifs, avec unité, et avec inverses \deleted{pour} là où il doit y

% archive sheet number: 384
